
\section{Limitations and Critical Reflection}
\label{sec:limitation}
Our method achieves large performance improvements on the PSG task, but several limitations and unexamined assumptions remain.

\noindent \textbf{The language information is static and pre-extracted.}
Our approach requires pre-extracting language information from LLMs before relation prediction. The extraction is efficient and cost-effective, but it creates a dependency on external models and on a preprocessing stage, which may limit scalability in some deployment settings. A more important consequence is that we elicit descriptions \emph{per category pair} rather than per image, so the descriptions are fixed at training time. This is what makes the method affordable, since the LLM is queried once for each subject-object category combination instead of once per test image. It also means that the language signal is a class-level prior rather than an instance-level observation. The model receives the same description of a person \textit{cleaning} an elephant regardless of the particular elephant in the image. When a specific instance differs from the typical case, the language branch cannot detect the difference, and the gating network can only discount the branch rather than correct it.

\noindent \textbf{Hallucination and the direction of the error.}
LLM-generated descriptions are not grounded in the image, and no component of our pipeline verifies them. An LLM asked how a person and a surfboard might be related produces fluent and plausible text whether or not that relation is present, and it describes the \emph{stereotypical} interaction because that is what its training distribution supports. The hallucination risk therefore has a specific and uncomfortable shape: the errors it introduces are biased toward the common case, which is the bias this chapter aims to reduce. The gating network is our safeguard, since it learns when to trust the language signal, but it is trained on the same long-tailed data, so its calibration in the tail is learned from few examples. Our ablation shows that GPT-4 and Llama3-8B descriptions perform comparably (Sec.~\ref{vlprompt:sec:experiments}), which is reassuring about robustness to the choice of LLM. It does not show that either set of descriptions is factually grounded, because two models that share a stereotype produce equally consistent and equally unfounded text. A stronger design would score descriptions against visual evidence before fusion, instead of leaving the entire decision to a learned gate.

\noindent \textbf{The improvements are measured against an incomplete annotation.}
As discussed in Sec.~\ref{hilo:sec:limitations}, R@$K$ and mR@$K$ score predictions against PSG annotations that are known to omit valid relations. The central result of this chapter, a 13.8 point mR@100 gain on rare relations, is therefore a claim about agreement with a particular annotation rather than about correctness in general. The concern is sharper here than in Chapter~3. Language priors could in principle improve the metric by aligning the output distribution with the linguistic conventions that annotators themselves used, without improving visual grounding. The predicate-wise analysis and the w/o-LLM ablation argue against this reading, because the gains concentrate in relations that are semantically distinguishable rather than spreading uniformly. The experiments in this chapter cannot rule it out entirely.

\noindent \textbf{Failure is concentrated in one diagnosable case.}
The failure analysis in Sec.~\ref{vlprompt:sec:experiments} shows that errors cluster where several instances of one category appear close together. The model then attaches a relation to the wrong instance, or the panoptic backbone merges two instances into a single mask and makes the correct assignment impossible. This is a scope limit worth stating plainly. Language priors operate at the level of \emph{categories}, describing how a person relates to a baseball glove, so they carry no information about \emph{which} person or \emph{which} glove. The mechanism contributed by this chapter therefore cannot resolve its own most frequent failure mode, and progress on that case requires better instance-level grounding rather than better priors.

\noindent \textbf{Closed-set by construction.}
Our method targets closed-set relation prediction and does not support open-set relation prediction, which restricts its use in real-world settings where unseen relations occur. This restriction is architectural rather than incidental, since descriptions are enumerated over a fixed set of category pairs and the prompter outputs a distribution over a fixed predicate list. We significantly reduce the long-tail problem within the closed-set setting, but closing the gap to open-set generalization remains an important challenge. Chapter~5 addresses this gap directly by extending PSG from closed-set relation classification to open-set relation prediction with large multimodal models.

\section{Conclusion}
\label{sec:conclusion}

This chapter showed that language priors from LLMs can complement the closed-set visual modeling established in Chapter~3.
VLPrompt incorporates language information generated by LLMs to enhance PSG performance.
VLPrompt utilizes the chain-of-thought method in designing prompts, enabling LLMs to generate rich descriptions for relation prediction. 
Additionally, we develop a prompter network based on attention mechanisms to facilitate comprehensive interaction between vision and language information, achieving high-quality relation prediction.
Experiments demonstrate that our method significantly outperforms the current state-of-the-art on the PSG dataset and mitigates the long-tail problem for relations.
Together, Chapters~3 and~4 improve PSG under long-tail and ambiguous supervision while preserving visual grounding.
The next chapter extends this line beyond a fixed relation vocabulary by studying open-set panoptic scene graph generation.
